%% LyX 2.4.4 created this file.  For more info, see https://www.lyx.org/.
%% Do not edit unless you really know what you are doing.
\documentclass[english]{article}
\usepackage[T1]{fontenc}
\usepackage[utf8]{inputenc}
\setcounter{secnumdepth}{-2}
\setlength{\parindent}{0bp}
\usepackage{amsmath}
\usepackage{amsthm}
\usepackage{amssymb}

\makeatletter
%%%%%%%%%%%%%%%%%%%%%%%%%%%%%% Textclass specific LaTeX commands.
\theoremstyle{definition}
\newtheorem*{problem*}{\protect\problemname}

\@ifundefined{date}{}{\date{}}
%%%%%%%%%%%%%%%%%%%%%%%%%%%%%% User specified LaTeX commands.
\usepackage{graphicx}
\usepackage{geometry}
\newcommand{\limi}{\underset{n\rightarrow\infty}{\lim}}
\newcommand{\limxz}{\underset{x\rightarrow x_{0}}{\lim}}
\newcommand{\limfxz}{\underset{x\rightarrow x_{0}}{\lim}f(x)}
\newcommand{\limfxm}{\underset{x\rightarrow x_{0}^{-}}{\lim}f(x)}
\newcommand{\limfxp}{\underset{x\rightarrow x_{0}^{+}}{\lim}f(x)}
\newcommand{\limxm}{\underset{x\rightarrow x_{0}^{-}}{\lim}}
\newcommand{\limxp}{\underset{x\rightarrow x_{0}^{+}}{\lim}}
\newcommand{\sqrm}{M_{m\times m}(\mathbb{F})}
\newcommand{\seqa}{(a_{n})_{n=1}^{\infty}}
\newcommand{\seqx}{(x_{n})_{n=1}^{\infty}}
\newcommand{\funcdr}{f:D\rightarrow\mathbb{R}}
\newcommand{\der}{\limxz\frac{f(x)-f(x_{0})}{x-x_{0}}}
\usepackage[all]{pdfprivacy}

\makeatother

\usepackage{babel}
\providecommand{\problemname}{Problem}

\begin{document}
\title{{\Large Semester 2026A, Exam A, Q9\vspace{-2em}}}
\maketitle
\begin{problem*}
Are the following matrices similar? 
\[
A=\begin{pmatrix}1 & -1 & 1 & -1 & 1\\
-1 & 1 & -1 & 1 & -1\\
1 & -1 & 1 & -1 & 1\\
-1 & 1 & -1 & 1 & -1\\
1 & -1 & 1 & -1 & 1
\end{pmatrix},\qquad B=\begin{pmatrix}1 & 0 & 1 & 0 & 1\\
0 & 1 & 0 & 1 & 0\\
1 & 0 & 1 & 0 & 1\\
0 & 1 & 0 & 1 & 0\\
1 & 0 & 1 & 0 & 1
\end{pmatrix}\;\in M_{5\times5}(\mathbb{R})
\]
\end{problem*}
\medskip{}

\rule[0.5ex]{1\columnwidth}{1pt}

\medskip{}

\begin{proof}
Let us recall the claim that if matrices are similar, they have the
same rank; thus, if matrices have different ranks, they are not similar.

\[
\text{rk}\text{A =\text{rk} }T_{A}=\text{dim}\left(Ax\mid x\in\mathbb{R}^{5}\right)
\]

\[
\text{rk}\text{B =\text{rk} }T_{B}=\text{dim}\left(Bx\mid x\in\mathbb{R}^{5}\right)
\]

For $x\in\mathbb{R}^{5}$,

\[
Ax=\begin{pmatrix}1 & -1 & 1 & -1 & 1\\
-1 & 1 & -1 & 1 & -1\\
1 & -1 & 1 & -1 & 1\\
-1 & 1 & -1 & 1 & -1\\
1 & -1 & 1 & -1 & 1
\end{pmatrix}\begin{pmatrix}x_{1}\\
x_{2}\\
x_{3}\\
x_{4}\\
x_{5}
\end{pmatrix}=\begin{pmatrix}x_{1}-x_{2}+x_{3}-x_{4}+x_{5}\\
-x_{1}+x_{2}-x_{3}+x_{4}-x_{5}\\
x_{1}-x_{2}+x_{3}-x_{4}+x_{5}\\
-x_{1}+x_{2}-x_{3}+x_{4}-x_{5}\\
x_{1}-x_{2}+x_{3}-x_{4}+x_{5}
\end{pmatrix}
\]

so 
\[
\text{Im}\left(T_{A}\right)=\left\{ t\begin{pmatrix}1\\
-1\\
1\\
-1\\
1
\end{pmatrix}\mid t=x_{1}-x_{2}+x_{3}-x_{4}+x_{5}\in\mathbb{R}\right\} 
\]

And since the scalar $t=x_{1}-x_{2}+x_{3}-x_{4}+x_{5}\in\mathbb{R}$
attains every value in $\mathbb{R}$,

\[
=\text{Span}\left(\begin{pmatrix}1\\
-1\\
1\\
-1\\
1
\end{pmatrix}\right)
\]

A sequence of one non zero vector is always linearly independent,
hence this sequence is a base of $\text{Im}\left(T_{A}\right)$ ,
so $\text{rk}A=1$.

For $x\in\mathbb{R}^{5}$,

\[
Bx=\begin{pmatrix}1 & 0 & 1 & 0 & 1\\
0 & 1 & 0 & 1 & 0\\
1 & 0 & 1 & 0 & 1\\
0 & 1 & 0 & 1 & 0\\
1 & 0 & 1 & 0 & 1
\end{pmatrix}\begin{pmatrix}x_{1}\\
x_{2}\\
x_{3}\\
x_{4}\\
x_{5}
\end{pmatrix}=\begin{pmatrix}x_{1}+x_{3}+x_{5}\\
x_{2}+x_{4}\\
x_{1}+x_{3}+x_{5}\\
x_{2}+x_{4}\\
x_{1}+x_{3}+x_{5}
\end{pmatrix}
\]

so 
\[
\text{Im}\left(T_{B}\right)=\left\{ t\begin{pmatrix}1\\
0\\
1\\
0\\
1
\end{pmatrix}+k\begin{pmatrix}0\\
1\\
0\\
1\\
0
\end{pmatrix}\mid t=x_{1}+x_{3}+x_{5}\in\mathbb{R},k=x_{2}+x_{4}\right\} 
\]

These scalars reach every value in $\mathbb{R}$ so

\[
=\text{Span }\left\{ \begin{pmatrix}1\\
0\\
1\\
0\\
1
\end{pmatrix},\begin{pmatrix}0\\
1\\
0\\
1\\
0
\end{pmatrix}\right\} 
\]

It's obvious this sequence is linearly independent (recall that a
sequence of two vectors is linearly dependent if and only if one vector
is a scalar multiple of the other):

\[
\forall x\in\mathbb{R}\quad x\begin{pmatrix}1\\
0\\
1\\
0\\
1
\end{pmatrix}=\begin{pmatrix}x\\
0\\
x\\
0\\
x
\end{pmatrix}\neq\begin{pmatrix}0\\
1\\
0\\
1\\
0
\end{pmatrix},\quad\wedge\quad\forall x\in\mathbb{R}\quad x\begin{pmatrix}0\\
1\\
0\\
1\\
0
\end{pmatrix}=\begin{pmatrix}0\\
x\\
0\\
x\\
0
\end{pmatrix}\neq\begin{pmatrix}1\\
0\\
1\\
0\\
1
\end{pmatrix}
\]

Therefore it's a basis for $\text{Im}\left(T_{B}\right)$, and $\text{rk}B=2\neq\text{rk}A$
which means these matrices are not similar.
\end{proof}
\begin{quote}
This is an exam problem I translated and solved. The original exam
was written by the Linear Algebra 1 course lecturers at HUJI.
\end{quote}

\end{document}
